\bookchapter{Appendice D: Dove trovare risorse affidabili?}

\booksection{Il supporto ufficiale è la prima scelta:}

\begin{itemize}
  \item \textbf{Utenti Apple}: apple.com/cn e iCloud.com.
  \item \textbf{Utenti Android}: sul sito del marchio del telefono c'è sempre una sezione "Servizio e supporto".
  \item \textbf{Strumenti universali}: Apple usa iTunes, Huawei usa Huawei Phone Assistant, Xiaomi usa Mi Mover.
\end{itemize}

\booksection{Siti e creator utili:}

\begin{itemize}
  \item \textbf{Trova trucchi}: cerca su Bilibili "[modello di telefono] trucchi".
  \item \textbf{Recensioni}: guarda le recensioni di più creator diversi.
  \item \textbf{Fotografia}: segui creator specializzati in tutorial di fotografia con telefono.
\end{itemize}

\booksection{Checklist di autoverifica sicurezza in una frase:}

\begin{enumerate}
  \item Hai impostato la password di blocco schermo?
  \item Hai attivato "Trova il mio telefono"?
  \item Hai attivato l'autenticazione a due fattori per gli account importanti?
  \item Fai backup periodici di foto e chat WeChat?
  \item Ricorda la formula anti-truffa: non credere ciecamente, non rivelare dati, non trasferire denaro, non essere avido, non essere precipitoso.
\end{enumerate}
